\section{Birational comparisons}
\label{sec:comparisons}

We compare the adjoints on the two sides of a small diagram and at the
ends of a finite relative program.  The comparisons locate exactly which
global places acquire larger discrepancies.  They also give a projective
morphism from a nef end to a prescribed ample model, with the nef adjoint
equal to the pullback of the ample adjoint.  We then record the homological
ranks that count contracted divisors and surfaces.
All divisor equalities in this section are equalities of actual rational
divisors, with the compatible canonical representatives and the fixed
b-divisor $\M$ of Section~\ref{sec:preliminaries}.

For a projective morphism, a rational Cartier divisor is called
\emph{relatively nef} here if it has nonnegative degree on every compact
irreducible curve contracted by the morphism.  We use the analytic
negativity lemma in the following precise form
\cite[Lemma~2.2]{CanonicalKahler}: if $r:V\to T$ is a projective
bimeromorphic morphism of normal complex spaces, $H$ is a
$\Q$-Cartier divisor on $V$, $-H$ is relatively nef, and $r_*H\geq0$,
then $H\geq0$.  Neither compactness nor a singularity hypothesis is
part of this input.  A divisor is \emph{exceptional over $T$} if every
prime in its support has image of codimension at least two in $T$.

We will also use two consequences of normality.  An effective
$\Q$-Cartier divisor pulls back effectively under a dominant morphism:
after clearing its index, a local meromorphic equation with no
divisorial poles is holomorphic on a normal space.  Such a divisor has
nonnegative degree on a compact curve not contained in its support,
by restricting its effective section to the normalization of the curve.

\subsection{Common models and the nef part}

The projectivity in the hypotheses permits all comparisons in a finite
part of the argument to be made on one compact K\"ahler resolution.
We give the construction for the diagrams used here; compare the
common-model construction of \cite[Proposition~2.1]{CanonicalKahler}.

\begin{lemma}[Common models with projective maps]\label{lem:common-projective-models}
Let a finite collection of normal irreducible compact K\"ahler spaces
be joined, with compatible bimeromorphic identifications, by finitely
many projective bimeromorphic morphisms and their inverses.  Suppose
that a fixed nef b-divisor is carried by a projective bimeromorphic
morphism $X'\to X$ to one of these spaces, with analytically nef
$\Q$-Cartier trace $M'$ on $X'$.  Then there is a
smooth compact K\"ahler space $W$ with compatible projective
bimeromorphic morphisms to every space in the collection and to $X'$.
The resolution may simultaneously principalize finitely many specified
generically nonzero coherent ideals and resolve finitely many specified divisors.  Its trace
$M_W$ is the pullback of $M'$ and is analytically nef.
\end{lemma}

\begin{proof}
For a diagram $T\to Z\leftarrow T^+$, take the reduced component of
$T\times_ZT^+$ containing the graph over the common isomorphism open.
Its projections are projective, and finite normalization preserves
projectivity.  Adjoin the remaining spaces and the carrier one at a
time by the same construction.  Resolving the resulting dominating
component gives a common smooth model.  Analytic resolution and
principalization allow the stated simultaneous choices by projective
modifications; see \cite[Theorems~1.6, 1.10 and 13.2(1)]{BierstoneMilman1997}.
Base change, restriction to a closed analytic subspace, finite
normalization, and composition preserve projectivity.  For composition
in the present compact setting, one can choose a uniform sufficiently
large twist of a relative ample bundle by the pullback of the next one.
This also proves compatibility with any finite chain of the small
diagrams and divisorial contractions considered below.

A smooth space projective over a compact K\"ahler space is compact
K\"ahler: a relative ample bundle supplies positivity in the fiber
directions, and adding a sufficiently large multiple of a pulled-back
K\"ahler form supplies positivity in all directions.  Finally, choose
K\"ahler classes on $X'$ converging to $c_1(M')$.  Their pullbacks to
$W$ have semipositive representatives.  Adding positive multiples
tending to zero of a fixed K\"ahler class on $W$ gives K\"ahler classes
converging to $c_1(M_W)$.  Thus $M_W$ is analytically nef.
\end{proof}

In particular, $M_W$ has nonnegative degree on every compact curve.
For an arbitrary global place, a further common proper model with a
model carrying that place suffices to compare coefficients.  This last
model need not be projective: projectivity is used to establish an
effective divisor on a common model, after which its pullbacks remain
effective on every higher proper model.

\begin{lemma}[Effective nef defect]\label{lem:nef-defect}
Let $T$ be a normal compact K\"ahler model, and let $r:W\to T$ be a
projective bimeromorphic morphism from a normal compact K\"ahler space
$W$ dominating the fixed projective carrier of $\M$.  Assume that
$M_T$ is $\Q$-Cartier.  Then
\begin{equation}\label{eq:nef-defect}
 J_{W/T}:=r^*M_T-M_W\geq0
\end{equation}
is an exceptional $\Q$-Cartier divisor over $T$.  If $s:W'\to W$ is
a higher normal model, then $J_{W'/T}=s^*J_{W/T}$.
\end{lemma}

\begin{proof}
The trace definition gives $r_*J_{W/T}=0$.  On an $r$-contracted curve,
$-J_{W/T}=M_W-r^*M_T$ has the same degree as $M_W$, hence has
nonnegative degree by projection formula from the nef carrier.
Analytic negativity gives $J_{W/T}\geq0$.
The final identity follows from functoriality of pullback and the fact
that the b-divisor descends on $W$.
\end{proof}

The next lemma explains why hypotheses stated for global generalized
places give the local ordinary singularity conditions needed later.
It is here that effectivity of the boundary and of the nef defect
enters the singularity comparison.  The global Weil and strong
$\Q$-factoriality conventions ensure the Cartier hypotheses below for
the respective models in the proof, since the divisor traces have
finite support.

\begin{lemma}[From global generalized discrepancies to local singularities]
\label{lem:ordinary-singularities}
Let $T$ be a normal compact K\"ahler space, let $B$ be an effective
rational divisor of finite support, and let $\M$ have analytically nef
rational data on a projective carrier.  Assume that $K_T$, $B$, and
$M_T$ are $\Q$-Cartier.
If every global generalized log discrepancy of $(T,B+\M)$ is
nonnegative, respectively positive, then the same condition holds for
local analytic places.  In the positive case, $T$ has ordinary klt
singularities locally.  If the global generalized discrepancies are
positive and every global place exceptional over $T$ has generalized
log discrepancy strictly greater than one,
then $T$ has ordinary terminal singularities locally.  In dimension
four the latter conclusion implies $\dim\Sing T\leq1$.
\end{lemma}

\begin{proof}
Choose a smooth log resolution $r:W\to T$, projective over $T$ and carrying $\M$,
with divisorial exceptional locus and simple normal crossing support
for all divisors in the crepant equation
\[
 K_W+\Delta_W+M_W=r^*(K_T+B+M_T).
\]
Such a simultaneous resolution is available by
Lemma~\ref{lem:common-projective-models}.  The assumed inequalities
give coefficients at most one, respectively strictly below one, for
every global prime of $\Delta_W$.  The same bounds hold on their local
branches.  On any further local model the nef part is pulled back from
$W$, so its generalized discrepancies are the ordinary discrepancies
of the simple normal crossing subpair $(W,\Delta_W)$.  The simple
normal crossing discrepancy criterion proves the asserted local
generalized inequalities, including when some coefficients of
$\Delta_W$ are negative.

Put $J=J_{W/T}$.  The ordinary crepant boundary for $(T,0)$ is
\begin{equation}\label{eq:ordinary-crepant-boundary}
 \Delta_W^{(0)}=\Delta_W-r^*B-J,
 \qquad K_W+\Delta_W^{(0)}=r^*K_T.
\end{equation}
Both $r^*B$ and $J$ are effective.  After resolving their joint support
as above, every coefficient of $\Delta_W^{(0)}$ is therefore strictly
below one in the positive case.  The same local simple normal crossing
criterion gives ordinary klt singularities.

For the terminal assertion, the relative canonical divisor is
exceptional, and at an exceptional prime $R$ of $W$ its coefficient is
\[
 \operatorname{coeff}_R(K_W-r^*K_T)
 =a(R;T,B+\M)-1+\operatorname{coeff}_R(r^*B+J)>0.
\]
These positive coefficients persist on every local branch.  A local
place exceptional over the smooth space $W$ has ordinary log
discrepancy at least two over $W$: the Jacobian determinant of a smooth
model carrying it has positive integral order along an exceptional
divisor.  Its discrepancy over $T$ adds the nonnegative order of
$K_W-r^*K_T$, and hence remains greater than one.  Local exceptional
places already represented by divisors on $W$ satisfy the displayed
strict inequality.  Thus $T$ is locally ordinary terminal.
For a normal complex fourfold with a rational canonical line bundle,
local ordinary terminality implies smoothness in codimension two by
\cite[Lemma~2.3]{CanonicalKahler}, giving the last assertion.
\end{proof}

For later calculations, the same equations give, for every global place
$E$ followed on a common higher model,
\begin{equation}\label{eq:generalized-ordinary-discrepancy}
 a(E;T,B+\M)
 =a(E;T,0)-\ord_E(r^*B)-\ord_E(J_{W/T}).
\end{equation}
The orders on the right mean coefficients after pullback to that model.
They are nonnegative under the hypotheses of
Lemma~\ref{lem:ordinary-singularities}.

\subsection{Small diagrams and finite programs}

We now locate the entire region in which a small step changes
discrepancies.  All exceptional loci and their images in this
subsection are given their reduced structures.  We use the graph method
of \cite[Lemma~2.1]{ProjectiveLCFourfolds} and prove the needed analytic
comparison, with the fixed b-divisor, below.

\begin{lemma}[Comparison for an arbitrary small diagram]
\label{lem:small-comparison}
Let
\[
 T\xrightarrow{\ f\ }Z\xleftarrow{\ f^+\ }T^+
\]
be a diagram of normal irreducible compact K\"ahler fourfolds in which
$f$ and $f^+$ are projective small bimeromorphic morphisms with
connected fibers, and both are nonisomorphisms.  Let $(T,B_T+\M)$ and
$(T^+,B_{T^+}+\M)$ have effective rational boundaries related by
strict transform and the same fixed rational b-divisor.  Transport the
canonical representatives by the identification of prime divisors,
and assume that the actual adjoints $D_T$ and $D_{T^+}$ are
$\Q$-Cartier, with $-D_T$ $f$-ample and $D_{T^+}$ $f^+$-ample.

The exceptional images of $f$ and $f^+$ are equal.  If their common
image is $A$, then
\begin{equation}\label{eq:small-loci}
 L:=\Exc(f)=f^{-1}(A),\qquad
 L^+:=\Exc(f^+)=(f^+)^{-1}(A),
\end{equation}
the complements are isomorphic, and
\begin{equation}\label{eq:small-dimensions}
 \dim L,\dim L^+\leq2,\qquad \dim A\leq1,
 \qquad \dim L+\dim L^+\geq3.
\end{equation}

Let $G$ be the normalization of the main graph, with projections
$p:G\to T$, $q:G\to T^+$ and $h=f\circ p=f^+\circ q$.  Then the
actual $\Q$-Cartier divisor
\begin{equation}\label{eq:small-divisor}
 F:=p^*D_T-q^*D_{T^+}
\end{equation}
is effective and exceptional over both $T$ and $T^+$, and
\begin{equation}\label{eq:small-support}
 \Supp F=h^{-1}(A).
\end{equation}
Moreover, $F\cdot C<0$ for every irreducible curve $C$ contracted by $h$.

For every global place $E$ over these spaces, let $\ord_E(F)$ denote
the coefficient of the pullback of $F$ on any higher proper model
dominating $G$ and carrying $E$.  Then
\begin{equation}\label{eq:small-discrepancies}
 a(E;T^+,B_{T^+}+\M)-a(E;T,B_T+\M)=\ord_E(F)\geq0.
\end{equation}
The inequality is strict exactly when $\cent_T(E)\subset L$, and this
condition is equivalent to $\cent_{T^+}(E)\subset L^+$.  A place is
exceptional over $T$ if and only if it is exceptional over $T^+$.
\end{lemma}

\begin{proof}
We first recall the fiber description for a proper bimeromorphic
morphism $v:V\to S$ of normal spaces.  Its positive-dimensional-fiber
locus is closed analytic.  Over its complement the map is finite
locally on $S$, hence an isomorphism by normality.  Stein factorization
shows that every fiber is connected, because its finite bimeromorphic
factor over the normal space $S$ is an isomorphism.  A reduced compact
connected positive-dimensional fiber has no zero-dimensional
irreducible component: such a component would be a singleton disjoint
from the finitely many other closed components.  Hence no point in
such a fiber is a local isomorphism point.  The exceptional locus of
$v$ is exactly the full inverse image of its positive-dimensional-fiber
locus.

Let $U\subset Z$ be the open set over which $f$ is an isomorphism.
On $U$, $D_T$ defines a $\Q$-Cartier divisor $D_U$.  The actual Weil
coefficients of the two adjoints agree under the identification of
prime divisors.  Since $f^+$ is small, it follows that
$D_{T^+}|_{(f^+)^{-1}U}=(f^+)^*D_U$.  A positive-dimensional projective
fiber over $U$ would contain a curve on which this divisor has degree
zero, contrary to $f^+$-ampleness.  Thus $f^+$ is an isomorphism over
$U$.  Interchanging the two sides and using $f$-ampleness of $-D_T$
gives the reverse inclusion of the isomorphism opens.  This proves
the equality of exceptional images and \eqref{eq:small-loci}.
Smallness gives $\dim L,\dim L^+\leq2$.  Over a general point of each
component of $A$ the fiber has dimension at least one, so the
fiber-dimension theorem gives $\dim A\leq1$.

The graph projections are projective and surjective, and the map from
$G$ to its component in $T\times_ZT^+$ is finite.  A prime divisor of
either side maps to a prime divisor of $Z$ by smallness.  On the graph
it therefore corresponds to the same prime place on the other side,
with the same adjoint coefficient.  Consequently $p_*F=q_*F=0$.
On a $p$-contracted curve, $-F$ has the degree of $q^*D_{T^+}$, which
is nonnegative.  Analytic negativity gives $F\geq0$.

For a curve $C$ contracted by $h$, the projection formula and the two
ample signs give
\[
 -F\cdot C=(-D_T)\cdot p_*C+D_{T^+}\cdot q_*C>0.
\]
Here an image that is a point contributes zero.  At least one image is
a curve, because the finite map to the fiber product cannot contract
$C$ to a point.  Thus at least one summand is strictly positive.

The support of $F$ is contained in $h^{-1}(A)$, since both projections
identify the common isomorphism open and the adjoints agree there.
Conversely, take $x\in h^{-1}(A)$ and write $z=h(x)$.  The map
$h:G\to Z$ is proper bimeromorphic, so its fiber over $z$ is connected.
It is projective and positive-dimensional: its projection to the
positive-dimensional fiber $f^{-1}(z)$ is surjective.  The fiber
description above places $x$ on a positive-dimensional irreducible
component.  Successive projective hyperplane cuts through $x$ give an
irreducible curve $C$ in that fiber through $x$.  If $x$ were outside
$\Supp F$, then $C$ would not be contained in the effective
$\Q$-Cartier divisor $F$, so $F\cdot C\geq0$.  The preceding strict
inequality excludes this.  Hence \eqref{eq:small-support} holds.

Follow a global place $E$ on a common proper model dominating $G$ and
the nef carrier.  Subtracting its two crepant equations cancels the
same trace of $\M$ and gives \eqref{eq:small-discrepancies}.  The order
of an effective $\Q$-Cartier divisor is positive precisely when the
center of the place is contained in its support.  Proper images of
centers give
\[
 \cent_Z(E)=f(\cent_T(E))=f^+(\cent_{T^+}(E))=h(\cent_G(E)).
\]
Together with \eqref{eq:small-loci} and \eqref{eq:small-support}, this
proves both center equivalence and the stated exact strictness
criterion.  If the centers are not contained in these loci, their
general points correspond on the common open and the order is zero.
Finally, a place centered at a divisor on a normal model is that
divisorial place.  Smallness identifies these prime divisors on the
two sides, proving equivalence of exceptionality.

Since both morphisms are nonisomorphisms, $A$ is nonempty, so $F$ has
a prime divisor $R$ in its support.  The graph is a fourfold, hence
$\dim R=3$.  Its finite map to the graph component has image in
$L\times_ZL^+$, whence
\[
 3\leq\dim(L\times_ZL^+)\leq\dim L+\dim L^+.
\]
This is the remaining inequality in \eqref{eq:small-dimensions}.
\end{proof}

The discrepancy statement preserves generalized lc and generalized
klt inequalities along a small step.  It also preserves generalized
terminality, because the same places are exceptional on both sides.
The isomorphism-open argument did not use nontriviality.  In the normal
irreducible compact K\"ahler fourfold setting of
Lemma~\ref{lem:small-comparison}, if an elementary small step is initially specified by a nonisomorphic
negative contraction, a small positive morphism, and the stated Cartier
and ample hypotheses, that argument forces the positive morphism to be
nonisomorphic as well.  Its fibers are connected by normality and Stein
factorization, so Lemma~\ref{lem:small-comparison} applies.
For a finite program we need, in addition, to retain the location of
the accumulated correction at its end.

\begin{lemma}[Comparison along a finite negative program]
\label{lem:program-comparison}
Let $Y_0\dashrightarrow\cdots\dashrightarrow Y_m$, with $m\geq0$, be
a finite sequence of normal irreducible compact K\"ahler fourfolds with
$\Q$-Cartier adjoints $D_j=K_{Y_j}+B_j+M_{Y_j}$.  Suppose the effective
rational boundaries are transported by pushforward, the canonical
representatives are compatible, and the same b-divisor with a projective
carrier is used throughout.  Assume
each step is either a small diagram satisfying
Lemma~\ref{lem:small-comparison}, or a projective bimeromorphic
divisorial contraction $v_j:Y_j\to Y_{j+1}$ such that
$D_{j+1}=(v_j)_*D_j$ and $-D_j$ is $v_j$-ample.  These steps extract
no prime divisors.  On a smooth common model $W$ with projective maps
to every $Y_j$ and to the carrier, write $r_j:W\to Y_j$.  Then
\begin{equation}\label{eq:program-comparison}
 r_0^*D_0-r_m^*D_m=E,\qquad E\geq0,\qquad (r_m)_*E=0.
\end{equation}
Every global discrepancy is nondecreasing along the run.  If a prime
divisor is contracted in a divisorial step, its discrepancy increases
strictly at that step, and its discrepancy at the end is strictly
larger than before that step.
\end{lemma}

\begin{proof}
For a divisorial step put
$E_j=D_j-v_j^*D_{j+1}$.  It is exceptional over $Y_{j+1}$ and has
strictly negative degree on every $v_j$-contracted curve.  Negativity
gives $E_j\geq0$.  Every point of an exceptional fiber lies on a curve
in that fiber, by the connected projective fiber argument in the proof
of Lemma~\ref{lem:small-comparison}.  The negative degree therefore
forces that point into $\Supp E_j$.  In particular the coefficient of
$E_j$ at each contracted prime is positive.  Subtracting crepant
equations gives nondecreasing discrepancies and this strict increase.
For a small step the same assertions, with an effective correction
exceptional over its output, follow from
Lemma~\ref{lem:small-comparison}.

Pull the corrections of all steps to $W$ and add them.  This gives the
effective divisor in \eqref{eq:program-comparison}.  Each prime in a
summand is exceptional over that step's output, and remains exceptional
over every later output.  Indeed, if that place became a prime divisor
on a later normal model, that prime would have a strict-transform prime
on each earlier model, since none of the intervening steps extracts divisors.
This contradicts its earlier exceptionality.  Thus the sum is
exceptional over $Y_m$.  Its nonnegative coefficients also show that
a strict increase from a contracted prime cannot be canceled later.
For $m=0$ the correction is zero and the assertions are immediate.
\end{proof}

\subsection{A nef end and a prescribed ample model}

An effective correction exceptional over the end is the form needed to
compare a finite program with a prescribed positive side of a diagram.
The next argument uses the effective-pushforward form of negativity,
since the two corrections need not have the same exceptional primes;
compare \cite[Lemma~2.5]{ProjectiveLCFourfolds} in the projective setting.

\begin{lemma}[Nef and ample models]\label{lem:nef-ample}
Let $Y_0$, $Y$, and $T_+$ be normal irreducible compact K\"ahler spaces
with projective bimeromorphic morphisms to a normal compact K\"ahler
space $S$.  On a smooth common model $W$ over $S$, projective over these
spaces, let $P_0$, $P$, and $P_+$ be the pullbacks of respective
$\Q$-Cartier divisors $D_0$, $D_Y$, and $D_+$.  Suppose that
\[
 P_0=P+E_Y=P_++E_+,
\]
where $E_Y\geq0$ is exceptional over $Y$ and $E_+\geq0$ is exceptional
over $T_+$.  If $D_Y$ is nef over $S$ and $D_+$ is ample over $S$,
then $P=P_+$ as actual divisors.  The induced bimeromorphic map is a
projective morphism $\varphi:Y\to T_+$ over $S$, and
$D_Y=\varphi^*D_+$.
\end{lemma}

\begin{proof}
Write $r:W\to Y$ and $t:W\to T_+$, and put
$H=P-P_+=E_+-E_Y$.  Then $r_*H=r_*E_+\geq0$.  On an $r$-contracted
curve, $-H$ has the degree of $P_+$, which is nonnegative because
the curve lies over a point of $S$.  Negativity gives $H\geq0$.
Likewise $t_*(-H)=t_*E_Y\geq0$, while $H$ is $t$-nef because $P$ comes
from an $S$-nef divisor.  Negativity applied to $-H$ gives $H\leq0$.
Thus $P=P_+$.

For a curve $C$ in an $r$-fiber, equality gives
$t^*D_+\cdot C=P\cdot C=0$.  The image $t(C)$ lies in an $S$-fiber,
where $D_+$ is ample, so $t(C)$ is a point.  Every $r$-fiber is
projective and connected.  If an irreducible component of one such
fiber had positive-dimensional image under $t$, a curve in that
projective image and projective hyperplane cuts of its inverse image
would supply a curve on which $t$ is nonconstant.  Thus each component
maps to a point, and connectedness makes those points equal.  The map
$t$ is constant on every $r$-fiber.

The reduced image of $(r,t):W\to Y\times_ST_+$ is a proper analytic
graph with one-point fibers over $Y$.  Its projection to $Y$ is
therefore finite and bimeromorphic, hence an isomorphism by normality
of $Y$.  This gives the morphism $\varphi$.  Its graph is a closed
subspace of $Y\times_ST_+$, whose projection to $T_+$ is projective
by base change from $Y\to S$; hence $\varphi$ is projective.  Finally,
pushing the equality $P=P_+$ to $Y$ gives
$D_Y=\varphi^*D_+$ as actual divisors.
\end{proof}

\subsection{Ranks of analytic cycle classes}

The small-diagram comparison forces a surface on at least one side of
every step.  The following ranks count the steps after the positive
side has no exceptional surface.  They also count divisorial steps in
the auxiliary programs.  Their role is the analytic counterpart of the
surface-cycle rank in the terminal fourfold argument of
\cite[Theorem~5-1-15 and Lemma~5-1-17]{KawamataMatsudaMatsuki1987}.

For a compact complex analytic space $V$ and an integer $k\geq0$, put
\[
 \begin{split}
 \mathcal C_k(V)&=\operatorname{span}_{\Q}
 \{[Q]\in H_{2k}(V,\Q):
       Q\subset V\text{ irreducible compact analytic},\ \dim Q=k\},\\
 c_k(V)&=\dim_{\Q}\mathcal C_k(V).
 \end{split}
\]
Fundamental classes use the complex orientation of the regular loci.
Finite compatible triangulations of compact analytic spaces make these
ranks finite.  If $V$ is compact K\"ahler, the class of every such $Q$
is nonzero, including when $V$ or $Q$ is singular: for a K\"ahler form
$\omega$,
\[
 \langle[\omega]^k,[Q]\rangle=\int_{Q_{\mathrm{reg}}}\omega^k>0.
\]
The same observation applies to a closed analytic subspace of a
compact K\"ahler space.  We use ordinary homology on compact spaces
and Borel--Moore homology on open complements.

\begin{lemma}[Loss of analytic cycle classes]\label{lem:cycle-loss}
Let $k\geq0$ be an integer, and let $X\dashrightarrow Y$ be a
bimeromorphic transformation of compact K\"ahler spaces represented by
a proper bimeromorphic diagram.
Suppose it identifies
\[
 U=X\setminus A\simeq Y\setminus B,
\]
where $A$ and $B$ are closed analytic subspaces and $\dim B<k$.
Then restriction to $U$ and strict transform induce a surjection
$\mathcal C_k(X)\to\mathcal C_k(Y)$.  Consequently
$c_k(X)\geq c_k(Y)$, with strict inequality if $A$ contains an
irreducible compact analytic subspace of dimension $k$.

In a sequence of bimeromorphic transformations of normal irreducible $n$-dimensional
compact K\"ahler spaces which extract no prime divisors, only finitely
many transformations can contract a prime divisor; a small
transformation preserves $c_{n-1}$.  In particular, for a small
fourfold diagram of Lemma~\ref{lem:small-comparison}, if $\dim L^+<2$,
then $c_2(T)\geq c_2(T^+)$, strictly if $L$ contains a surface.
\end{lemma}

\begin{proof}
This is the rational-coefficient version of
\cite[Lemma~4.1]{KahlerAuxiliary}; we include its localization argument.
The Borel--Moore localization sequence for $B\subset Y$ contains
\[
 H_{2k}(B,\Q)\longrightarrow H_{2k}(Y,\Q)
 \longrightarrow H_{2k}^{\BM}(U,\Q)
 \longrightarrow H_{2k-1}(B,\Q).
\]
The outside groups vanish by $\dim B<k$ and the real dimension bound
for an analytic triangulation.  Thus restriction from $Y$ is an
isomorphism in degree $2k$.  Compose restriction from $X$ with its
inverse.  An irreducible compact analytic $k$-cycle meeting $U$ is
sent to its strict transform on $Y$; analyticity of that transform
follows by proper mapping through the given diagram.  A cycle
contained in $A$ is sent to zero.  Conversely, no $k$-cycle of $Y$ is
contained in $B$, and its strict transform on $X$ maps to its class.
This proves the surjection on cycle spans.  If $A$ contains a
$k$-dimensional irreducible analytic subspace, its nonzero class is in
the kernel, giving strict inequality.

For a transformation that extracts no prime divisors, choose the
common isomorphism open with target complement of codimension at least
two.  Every contracted source prime lies in the source complement.
The preceding result with $k=n-1$ gives a strict drop at each such
contraction.  The ranks are nonnegative integers, so there can be
only finitely many drops.  For a small transformation both complements
have codimension at least two, and applying the result in both
directions gives equality.  Finally, \eqref{eq:small-loci} gives the
common open for a small fourfold diagram.  Applying the first assertion
with $k=2$, $A=L$ and $B=L^+$ proves the last claim.
\end{proof}
